\paragraph{Thermodynamic Completion of the Energy Spectrum: Large Deviations, Microcanonical Entropy, and Multifractal Readout}\label{par:pom-multiplicity-energy-ldp-multifractal}
Continuing with the ordered composition set \(\mathcal{C}_L\), the weight
\(
w(\boldsymbol{\ell})=\prod_i F_{\ell_i+2}
\)
from Definition \ref{def:pom-multiplicity-composition-partition}, and the fiber weight measure
\(
P_L=\mathbb{P}_L
\)
from Definition \ref{def:pom-multiplicity-induced-composition-measure}.
Denote the energy density of the logarithmic multiplicity by
\[
E_L(\boldsymbol{\ell}):=\frac{1}{L}\log w(\boldsymbol{\ell})
\qquad(\boldsymbol{\ell}\in\mathcal{C}_L).
\]
By the endpoint bounds \(F_{L+2}\le w(\boldsymbol{\ell})\le 2^L\) from Proposition \ref{prop:pom-path-component-multiplicity-refinement-monotone-extrema} and \(L^{-1}\log F_{L+2}\to\log\varphi\), the reachable energy closure is
\[
E_L(\boldsymbol{\ell})\in\bigl[L^{-1}\log F_{L+2},\,\log 2\bigr],
\qquad
\mathrm{cl}\,\mathrm{Range}(E_L)\subseteq[\log\varphi,\log 2].
\]

\begin{theorem}[Exact large deviation principle for the logarithmic multiplicity under the fiber weight measure: Legendre spectrum of the pressure function \texorpdfstring{$\log\lambda(q)$}{log lambda(q)}]\label{thm:pom-multiplicity-energy-ldp-under-PL}
Let \(E_L\) be as above, and let \(\lambda(q)\) be the real-parameter moment pressure from Definition \ref{def:pom-multiplicity-real-q-pressure}.
Then the sequence of random variables \((E_L)_{L\ge 1}\) satisfies a large deviation principle under \((P_L)_{L\ge 1}\); the normalized cumulant generating function exists for \(t>-1\) and admits the closed form
\[
\boxed{\;
\Psi(t):=\lim_{L\to\infty}\frac{1}{L}\log\mathbb{E}_{P_L}\!\left[e^{tLE_L}\right]
=\log\lambda(1+t)-\log\lambda(1)\; }.
\]
The corresponding rate function is
\[
\boxed{\;
I(e)=\sup_{t>-1}\bigl\{t e-\Psi(t)\bigr\}\in[0,\infty]\; } ,
\]
and \(I\) is strictly convex and real analytic on the interior interval \((\log\varphi,\log 2)\), taking \(I(e)=+\infty\) outside this interval.
\par
Equivalently, taking the parameter \(q>0\) and setting
\[
e(q):=\frac{d}{dq}\log\lambda(q)
\quad(\text{which exists and has an explicit expression given by Proposition \ref{prop:pom-multiplicity-lambdaq-derivative-gibbs}}),
\]
one obtains the parametric closed form on the interior of the range:
\[
\boxed{\;
I\bigl(e(q)\bigr)=(q-1)e(q)-\log\frac{\lambda(q)}{\lambda(1)}\; } .
\]
\end{theorem}
\begin{proof}[Proof outline]
For any \(t>-1\) (writing \(q=1+t>0\)), we have
\[
\mathbb{E}_{P_L}\!\left[e^{tLE_L}\right]
=\sum_{\boldsymbol{\ell}\in\mathcal{C}_L}\frac{w(\boldsymbol{\ell})}{Z_L^{(1)}}\,w(\boldsymbol{\ell})^{t}
=\frac{Z_L^{(q)}}{Z_L^{(1)}},
\qquad
Z_L^{(q)}:=\sum_{\boldsymbol{\ell}\in\mathcal{C}_L}w(\boldsymbol{\ell})^{q}.
\]
By the concatenation structure and the generating function identity \( \sum_{L\ge 0}Z_L^{(q)}t^L=1/(1-W_q(t))\) (see \eqref{eq:pom-multiplicity-composition-moment-hierarchy-gf}, whose derivation holds verbatim for real \(q>0\)), combined with \(W_q(\rho(q))=1\) and \(\partial_tW_q(\rho(q))>0\), we obtain
\(
\lim_{L\to\infty}L^{-1}\log Z_L^{(q)}=\log\lambda(q)
\).
Hence \(\Psi(t)=\log\lambda(1+t)-\log\lambda(1)\).
\par
Since Theorem \ref{thm:pom-multiplicity-real-q-pressure} establishes that \(q\mapsto\log\lambda(q)\) is real analytic and strictly convex on \((0,\infty)\), the function \(\Psi\) is differentiable on \(t>-1\) and analytic in a neighborhood of \(t=0\); the G\"artner--Ellis theorem then yields the large deviation principle for \(E_L\) and the Legendre--Fenchel rate function.
The parametric form follows from \(e(q)=\Psi'(q-1)\) and the evaluation of \(I(e)=t e-\Psi(t)\) at the stationary point. \qed
\end{proof}

\begin{corollary}[Exponential rarity constants at boundary energies: \texorpdfstring{$\log(\lambda(1)/2)$}{log(lambda(1)/2)} and \texorpdfstring{$\log(\lambda(1)/\varphi)$}{log(lambda(1)/phi)}]\label{cor:pom-multiplicity-energy-ldp-boundary}
Let \(\varphi=\frac{1+\sqrt5}{2}\), and note that by the integer consistency of Theorem \ref{thm:pom-multiplicity-real-q-pressure},
\(
\lambda(1)=\lambda_1=\frac{3+\sqrt{17}}{2}
\) (see Proposition \ref{prop:pom-multiplicity-composition-partition-rational}).
Then the rate function of Theorem \ref{thm:pom-multiplicity-energy-ldp-under-PL} has finite boundary values at the endpoints of the reachable energy closure, satisfying
\[
\boxed{\;
\lim_{e\uparrow \log 2} I(e)=\log\frac{\lambda(1)}{2},\qquad
\lim_{e\downarrow \log\varphi} I(e)=\log\frac{\lambda(1)}{\varphi}\; } .
\]
Therefore, the ``fully fragmented state'' (\(\boldsymbol{\ell}=\boldsymbol{1}*L\)) and the ``single-block state'' (\(\boldsymbol{\ell}=(L)\)) are both linearly exponentially rare under \(P_L\), with exponential constants \(\log(\lambda(1)/2)\) and \(\log(\lambda(1)/\varphi)\) respectively.
\end{corollary}
\begin{proof}
From the weight \(w(\boldsymbol{1}*L)=2^L\) and the definition \(P_L(\boldsymbol{\ell})=w(\boldsymbol{\ell})/Z_L^{(1)}\), we obtain
\(
P_L(\boldsymbol{1}*L)=2^L/Z_L^{(1)}
\), hence
\[
-\lim_{L\to\infty}\frac1L\log P_L(\boldsymbol{1}*L)=\log\lambda(1)-\log 2=\log\frac{\lambda(1)}{2}.
\]
Similarly, for \(\boldsymbol{\ell}=(L)\) we have \(w(\boldsymbol{\ell})=F_{L+2}\) and \(L^{-1}\log F_{L+2}\to\log\varphi\), whence
\[
-\lim_{L\to\infty}\frac1L\log P_L(\boldsymbol{\ell}=(L))=\log\lambda(1)-\log\varphi=\log\frac{\lambda(1)}{\varphi}.
\]
Since both endpoints are realized by single atoms, the boundary values can also be viewed as endpoint extensions of \(I\); this is compatible with the well-behaved rate function structure of Theorem \ref{thm:pom-multiplicity-energy-ldp-under-PL}. \qed
\end{proof}

\begin{theorem}[Existence and complete closed-form duality of the microcanonical counting entropy]\label{thm:pom-multiplicity-microcanonical-entropy}
For \(e\in\RR\) and \(\delta>0\), define the energy shell count
\[
N_L(e,\delta):=\#\Bigl\{\boldsymbol{\ell}\in\mathcal{C}_L:\ |E_L(\boldsymbol{\ell})-e|<\delta\Bigr\}.
\]
Then the limit
\[
s(e):=\lim_{\delta\downarrow 0}\ \lim_{L\to\infty}\frac{1}{L}\log N_L(e,\delta)
\]
exists on \(e\in(\log\varphi,\log 2)\) and satisfies the thermodynamic duality identity
\[
\boxed{\;
\log\lambda(q)=\sup_{e\in[\log\varphi,\log 2]}\bigl(s(e)+q e\bigr)\qquad(q\ge 0)\; } ,
\]
whence
\[
\boxed{\;
s(e)=\inf_{q\ge 0}\bigl(\log\lambda(q)-q e\bigr)\qquad(e\in[\log\varphi,\log 2])\; } .
\]
Moreover, \(s\) is strictly concave and real analytic on the interior interval, with boundary behavior
\[
\boxed{\;
\lim_{e\downarrow \log\varphi}s(e)=\lim_{e\uparrow \log 2}s(e)=0\; } .
\]
\end{theorem}
\begin{proof}[Proof outline]
From the definition \(Z_L^{(q)}=\sum_{\boldsymbol{\ell}\in\mathcal{C}_L}e^{qLE_L(\boldsymbol{\ell})}\) and the limit
\(
\lim_{L\to\infty}L^{-1}\log Z_L^{(q)}=\log\lambda(q)
\)
(same as in the proof outline of Theorem \ref{thm:pom-multiplicity-energy-ldp-under-PL}), applying a Varadhan-type maximum term principle to the counting measure of \(E_L\) yields
\(
\log\lambda(q)=\sup_e(s(e)+q e)
\).
Taking the Legendre--Fenchel dual in \(q\) gives \(s(e)=\inf_{q\ge 0}(\log\lambda(q)-q e)\).
Strict concavity and analyticity follow by duality transfer from the strict convex analyticity of \(q\mapsto\log\lambda(q)\) (Theorem \ref{thm:pom-multiplicity-real-q-pressure}).
The endpoint \(s\to 0\) follows from the fact that the endpoint energies are each realized by a single composition (Proposition \ref{prop:pom-path-component-multiplicity-refinement-monotone-extrema}). \qed
\end{proof}

\begin{corollary}[Gibbs identity: exact relation between the energy large deviation under the fiber weight measure and the microcanonical entropy]\label{cor:pom-multiplicity-gibbs-identity}
Let \(I\) be the rate function of Theorem \ref{thm:pom-multiplicity-energy-ldp-under-PL}, and let \(s\) be the microcanonical counting entropy of Theorem \ref{thm:pom-multiplicity-microcanonical-entropy}.
Then on the interior interval \((\log\varphi,\log 2)\) the strict identity holds:
\[
\boxed{\;
I(e)=\log\lambda(1)-s(e)-e\; } ,
\]
which extends continuously to the endpoints (corresponding to the boundary values of Corollary \ref{cor:pom-multiplicity-energy-ldp-boundary}).
In particular, the typical energy
\(
e_\star=\left.\frac{d}{dq}\log\lambda(q)\right|_{q=1}
\)
uniquely satisfies the supporting line condition
\(
s'(e_\star)=-1
\), and
\(
s(e_\star)=\log\lambda(1)-e_\star
\).
\end{corollary}
\begin{proof}
From \(P_L(\boldsymbol{\ell})=e^{LE_L(\boldsymbol{\ell})}/Z_L^{(1)}\), the energy shell probability admits the approximate decomposition
\[
P_L\bigl(|E_L-e|<\delta\bigr)
=\sum_{|E_L-e|<\delta}\frac{e^{LE_L}}{Z_L^{(1)}}
=\frac{e^{Le+o(L)}}{Z_L^{(1)}}\,N_L(e,\delta).
\]
Taking logarithms, dividing by \(L\), and using \(L^{-1}\log Z_L^{(1)}\to\log\lambda(1)\) and \(L^{-1}\log N_L(e,\delta)\to s(e)\) yields
\(
I(e)=\log\lambda(1)-s(e)-e
\).
The supporting line condition follows from the stationary point equation of \(s(e)=\inf_q(\log\lambda(q)-q e)\). \qed
\end{proof}

\begin{proposition}[Multifractal spectrum of atomic weights: \texorpdfstring{$\alpha$}{alpha} interval and spectral function closed form]\label{prop:pom-multiplicity-atom-multifractal}
Let the single-atom negative log-probability density be
\[
\alpha_L(\boldsymbol{\ell}):=-\frac{1}{L}\log P_L(\boldsymbol{\ell})\qquad(\boldsymbol{\ell}\in\mathcal{C}_L).
\]
Then there is a uniform asymptotic
\[
\alpha_L(\boldsymbol{\ell})=\log\lambda(1)-E_L(\boldsymbol{\ell})+o(1),
\]
whence the reachable limit interval is
\[
\boxed{\;
\alpha\in\Bigl[\log\frac{\lambda(1)}{2},\ \log\frac{\lambda(1)}{\varphi}\Bigr]\; } .
\]
Moreover, for any interior \(\alpha\) (setting \(e=\log\lambda(1)-\alpha\)), the atomic exponent spectrum satisfies
\[
\boxed{\;
\lim_{\delta\downarrow 0}\ \lim_{L\to\infty}\frac{1}{L}\log
\#\Bigl\{\boldsymbol{\ell}\in\mathcal{C}_L:\ |\alpha_L(\boldsymbol{\ell})-\alpha|<\delta\Bigr\}
=f(\alpha):=s\bigl(\log\lambda(1)-\alpha\bigr)\; } ,
\]
where \(s\) is the microcanonical counting entropy of Theorem \ref{thm:pom-multiplicity-microcanonical-entropy}.
Therefore \(f\) is strictly concave and real analytic on the interior, with endpoint values \(f(\log(\lambda(1)/2))=f(\log(\lambda(1)/\varphi))=0\).
\end{proposition}
\begin{proof}[Proof outline]
From \(P_L(\boldsymbol{\ell})=w(\boldsymbol{\ell})/Z_L^{(1)}\) and \(L^{-1}\log Z_L^{(1)}\to\log\lambda(1)\), we obtain \(\alpha_L=\log\lambda(1)-E_L+o(1)\).
The counting spectrum conclusion follows from the bijection \(\alpha\leftrightarrow e\) combined with the \(s(e)\) result of Theorem \ref{thm:pom-multiplicity-microcanonical-entropy}. \qed
\end{proof}

\begin{proposition}[Continuous R\'enyi entropy rate/dimension spectrum: full real-number extension of integer results]\label{prop:pom-multiplicity-renyi-continuous}
For real \(q>0,q\ne 1\), define the R\'enyi entropy
\[
H_q(P_L):=\frac{1}{1-q}\log_2\sum_{\boldsymbol{\ell}\in\mathcal{C}_L}P_L(\boldsymbol{\ell})^{q}.
\]
Then the limit entropy rate exists and satisfies the closed form
\[
\boxed{\;
\lim_{L\to\infty}\frac{1}{L}H_q(P_L)
=\frac{1}{1-q}\Bigl(\log_2\lambda(q)-q\log_2\lambda(1)\Bigr)\; } .
\]
Moreover, the Shannon entropy rate is given by the differentiable limit:
\[
\boxed{\;
\lim_{L\to\infty}\frac{1}{L}H(P_L)
=\log_2\lambda(1)-\left.\frac{d}{dq}\log_2\lambda(q)\right|_{q=1}\; } .
\]
\end{proposition}
\begin{proof}
From \(P_L(\boldsymbol{\ell})=w(\boldsymbol{\ell})/Z_L^{(1)}\) we get
\(
\sum P_L^q=Z_L^{(q)}/(Z_L^{(1)})^q
\).
Taking logarithms, dividing by \(L\), and using \(L^{-1}\log Z_L^{(q)}\to\log\lambda(q)\) yields the R\'enyi entropy rate closed form; letting \(q\to 1\) and using differentiability of \(\log\lambda\) gives the Shannon entropy rate formula. \qed
\end{proof}

\begin{theorem}[Central limit theorem for the logarithmic multiplicity and the variance constant: second derivative of \texorpdfstring{$\log\lambda$}{log lambda}]\label{thm:pom-multiplicity-logw-clt}
Viewing \(\log w(\boldsymbol{\ell})\) as a random variable under \(P_L\), there exist constants
\[
e_\star:=\left.\frac{d}{dq}\log\lambda(q)\right|_{q=1},
\qquad
\sigma_\star^2:=\left.\frac{d^2}{dq^2}\log\lambda(q)\right|_{q=1}\in(0,\infty),
\]
such that the central limit theorem holds:
\[
\boxed{\;
\frac{\log w(\boldsymbol{\ell})-e_\star L}{\sqrt{L}}\ \Longrightarrow\ \mathcal{N}(0,\sigma_\star^2)\; } .
\]
Equivalently,
\(
\sqrt{L}\,(E_L-e_\star)\Longrightarrow \mathcal{N}(0,\sigma_\star^2)
\).
\end{theorem}
\begin{proof}[Proof outline]
By Theorem \ref{thm:pom-multiplicity-energy-ldp-under-PL}, \(\Psi(t)=\log\lambda(1+t)-\log\lambda(1)\) is analytic and strictly convex in a neighborhood of \(t=0\);
therefore a standard quasi-power/analytic perturbation central limit theorem can be applied (or alternatively, a Cram\'er-type second-order expansion of \(\Psi\) followed by Fourier inversion) to obtain the CLT for \(\log w\).
Positivity of the variance follows from strict convexity. \qed
\end{proof}

\begin{proposition}[Moderate deviation principle: quadratic rate fully determined by \texorpdfstring{$\sigma_\star^2$}{sigma^2}]\label{prop:pom-multiplicity-logw-mdp}
Take any sequence \(a_L\to\infty\) with \(a_L=o(\sqrt{L})\).
Then the moderate deviation principle holds under \(P_L\): for any \(x\in\RR\),
\[
\boxed{\;
\lim_{L\to\infty}\frac{1}{a_L^2}\log
\mathbb{P}\!\left(\frac{\log w(\boldsymbol{\ell})-e_\star L}{a_L\sqrt{L}}\ge x\right)
=-\frac{x^2}{2\sigma_\star^2}\; } ,
\]
and the same quadratic rate holds for the left tail.
\end{proposition}
\begin{proof}[Proof outline]
From the analyticity of \(\Psi\) near zero and its second-order expansion \(\Psi(t)=e_\star t+\frac12\sigma_\star^2 t^2+O(t^3)\), applying the moderate deviation version of the G\"artner--Ellis theorem with the scaling \(t=\theta\,a_L/\sqrt{L}\) yields the conclusion. \qed
\end{proof}

\begin{theorem}[Energy spectrum and rate under the uniform composition measure: Legendre dual of \texorpdfstring{$\log\lambda(q)-\log 2$}{log lambda(q)-log 2}]\label{thm:pom-multiplicity-energy-ldp-uniform}
Let \(U_L\) be the uniform measure on \(\mathcal{C}_L\) (so \(U_L(\boldsymbol{\ell})=2^{-(L-1)}\)).
Then \((E_L)_{L\ge 1}\) satisfies a large deviation principle under \((U_L)\), with cumulant generating function
\[
\boxed{\;
\Phi(q):=\lim_{L\to\infty}\frac{1}{L}\log\mathbb{E}_{U_L}\!\left[e^{qLE_L}\right]
=\log\lambda(q)-\log 2\qquad(q\ge 0)\; } ,
\]
and rate function
\[
\boxed{\;
J(e)=\sup_{q\ge 0}\bigl\{q e-\Phi(q)\bigr\}\; } .
\]
Furthermore, the endpoints satisfy
\[
\boxed{\;
\lim_{e\uparrow\log 2}J(e)=\lim_{e\downarrow\log\varphi}J(e)=\log 2\; } ,
\]
corresponding to the counting rarity of single-atom realizations at both ends (\(|\mathcal{C}_L|=2^{L-1}\)).
\end{theorem}
\begin{proof}[Proof outline]
From \(\mathbb{E}_{U_L}[e^{qLE_L}]=2^{-(L-1)}Z_L^{(q)}\) and \(L^{-1}\log Z_L^{(q)}\to\log\lambda(q)\) we get \(\Phi(q)=\log\lambda(q)-\log 2\).
Applying the G\"artner--Ellis theorem using differentiability and strict convexity of \(\Phi\) yields the LDP and the Legendre form of \(J\).
The endpoint values follow from the fact that the endpoint events are single atoms with uniform probability \(2^{-(L-1)}\). \qed
\end{proof}

\begin{theorem}[Canonical--microcanonical equivalence: local limit of the conditional energy shell as a \texorpdfstring{$q$}{q}-tilted i.i.d.\ process]\label{thm:pom-multiplicity-ensemble-equivalence}
For any \(q>0\), define the \(q\)-tilted fiber weight measure \(P_L^{(q)}\) (see Proposition \ref{prop:pom-multiplicity-low-temperature-gibbs}):
\[
P_L^{(q)}(\boldsymbol{\ell})=\frac{w(\boldsymbol{\ell})^{q}}{Z_L^{(q)}},
\qquad
Z_L^{(q)}=\sum_{\boldsymbol{\ell}\in\mathcal{C}_L}w(\boldsymbol{\ell})^{q}.
\]
Let \(p_q(k)=F_{k+2}^{\,q}\rho(q)^k\) (\(k\ge 1\)) be the i.i.d.\ part-length distribution induced by \(\rho(q)=\lambda(q)^{-1}\) (see \eqref{eq:pom-multiplicity-low-temp-iid-law}).
Write \(e(q)=\frac{d}{dq}\log\lambda(q)\).
\par
Then for any fixed window \(m\ge 1\) and any bounded continuous function \(\mathcal{F}\) (depending on the first \(m\) parts), the Gibbs conditioning principle holds:
\[
\boxed{\;
\lim_{\delta\downarrow 0}\ \lim_{L\to\infty}
\mathbb{E}_{P_L}\!\left[\mathcal{F}(\ell_1,\dots,\ell_m)\ \middle|\ |E_L-e(q)|<\delta\right]
=\mathbb{E}_{p_q^{\otimes m}}[\mathcal{F}]\; } .
\]
Therefore the microcanonical restriction on the energy density is locally equivalent to the \(q\)-tilted canonical ensemble; the special case \(q=1\) yields the local i.i.d.\ limit law for the fiber weight measure \(P_L\).
\end{theorem}
\begin{proof}[Proof outline]
Since \(P_L=P_L^{(1)}\) and
\(
\frac{dP_L^{(q)}}{dP_L}(\boldsymbol{\ell})\propto e^{(q-1)LE_L(\boldsymbol{\ell})}
\),
\(P_L^{(q)}\) is an exponential tilt of \(P_L\).
Applying the standard Gibbs conditioning principle to the rare event \(|E_L-e(q)|<\delta\) under \(P_L\) (guaranteed by the unique tilting parameter from the strictly convex rate function of Theorem \ref{thm:pom-multiplicity-energy-ldp-under-PL}), the conditional finite-window distribution asymptotically coincides with the finite-window distribution under \(P_L^{(q)}\).
By Proposition \ref{prop:pom-multiplicity-low-temperature-gibbs}, the finite-window limit under \(P_L^{(q)}\) is \(p_q^{\otimes m}\); combining yields the conclusion. \qed
\end{proof}

\endinput
